% /chapters/0-preamble/0-abstract.tex

\begin{abstract}
    A wide range of problems in computer science, including constraint satisfaction, can be framed in the language of graph homomorphisms.
    There are problems for which the existence of a solution may be proved by means of the Axiom of Choice or a non-principal ultrafilter.
    Although such solutions exist abstractly, in practice we want solutions that are definable, and so it is natural to impose restrictions.
    Interestingly, this often makes the problems harder to solve, but the solutions have a concrete description.
    This thesis studies the cost of insisting on definable solutions.
    The parallel between computer-science and set-theoretic complexity is a recurring theme in the thesis.

    Chapter~2 presents new proofs and results about the complexity of colouring problems for analytic graphs, with the main highlight being a resolution of the descriptive complexity of the 2-Borel-dicolourability problem, where we show that it is $\Sigma^1_2$-complete.
    This rules out a countable basis for Borel digraphs of dichromatic number at least three.
    We also present a streamlined proof of the Carroy--Miller--Schrittesser--Vidny{\'{a}}nszky $\mathbb{L}_0$ dichotomy, and the solution to an open problem about infinite games played on a linear order.

    In Chapter~3, we study the asymptotic behaviour of computations, motivated by neural networks of increasing depth, using tools from the topology of function spaces and, more specifically, the classification of Rosenthal compacta.
    We give explicit descriptions of new computational definability classes for the Alva--Due{\~n}ez--Iovino--Walton floating-point arithmetic framework, and establish connections to statistical learning theory.

    Finally, Chapter~4 studies a group action on graphs introduced by Caro, Hansberg, and Montejano to study balanceability and interpolation properties in extremal combinatorics, and develops a descriptive-set-theoretic theory of these actions on infinite graphs.

    Descriptive set theory provides the language for all three substantive chapters: Borel definability for Chapter~2, the Baire-class hierarchy for Chapter~3, and Polish group actions for Chapter~4.
\end{abstract}
